% Abhakseite „Das kann ich“ – Zone nur im Gesamt (\mitzone)
%% TODO Ich-kann-Sätze: Platzhalter wie in den Aufgabendateien
\begin{abhakseite}
\ifdefined\mitzone
\abhakgruppe{Kennst du schon}
\abhak{1}{Ableitungen bilden}
\abhak{2}{Ableitungswert als Tangentensteigung deuten (Grundvorstellung)}
\abhak{3}{Geradengleichung y = mx + n aus Steigung und Punkt aufstellen, Steigungsdreieck lesen und zeichnen, Nullstelle und y-Achsenabschnitt einer Geraden bestimmen, Punktprobe und Funktionswert}
\abhak{4}{Gleichungen lösen}
\abhak{5}{Tangens im rechtwinkligen Dreieck und Umkehrfunktion tan⁻¹, Grad-Modus des Rechners}
\abhak{6}{Satz des Pythagoras und Flächenformel des Dreiecks (halbes Produkt der Katheten), Streckenlänge zwischen zwei Punkten}
\abhak{7}{Ableitungswert als Tangentensteigung deuten (Grundvorstellung) – Fehler finden}
\abhak{8}{Ableitungswert als Tangentensteigung deuten (Grundvorstellung) – selbst rechnen}
\fi
\abhakgruppe{Einheit 1 · Tangentengleichung im Punkt}
\abhak{9}{Tangentengleichung}
\abhak{10}{Stellen, deren Tangente die y-Achse oberhalb einer Schranke schneidet, über eine Ungleichung bestimmen}
\abhak{11}{Parallelität der Wendetangenten einer Schar über eine parameterunabhängige Steigung nachweisen}
\abhak{12}{Tangentengleichung – Fehler finden, Begründen}
\abhakgruppe{Einheit 2 · Tangente als Berührung}
\abhak{13}{Berührung nachweisen}
\abhak{14}{Berührpunkt einer Tangente durch einen vorgegebenen Punkt berechnen}
\abhak{15}{Zeitpunkt bei gleichbleibender Änderungsrate über den Schnitt der Tangente mit der Zeitachse grafisch bestimmen}
\abhak{16}{Lage zweier Graphen mit gemeinsamen Endpunkten über die Steigungen in den Endpunkten begründen}
\abhak{17}{Berührung nachweisen – Fehler finden, Begründen}
\abhakgruppe{Einheit 3 · Normale}
\abhak{18}{Normale}
\abhak{19}{Normale – Fehler finden, Begründen}
\abhakgruppe{Einheit 4 · Steigungswinkel und Schnittwinkel}
\abhak{20}{Winkel}
\abhak{21}{Steigungswinkel einer Tangente als größten Kippwinkel im Sachzusammenhang deuten}
\abhak{22}{Winkel – Fehler finden, Begründen}
\abhakgruppe{Einheit 5 · Dreiecke und Figuren aus Tangente, Normale und Achsen}
\abhak{23}{Dreiecke und Figuren}
\abhak{24}{Dreiecke und Figuren – Fehler finden, Begründen}
\end{abhakseite}
